\chapter{Introdução}\label{cap:introducao}
%Trazer aplicações de RRT na prática, acho que libs como opml usam rrt e sao empregadas em roboticaespacial

A robótica, ao longo das últimas décadas, deixou de estar associada apenas à execução de tarefas repetitivas em ambientes rigidamente controlados, passando a representar uma área central no desenvolvimento de sistemas autônomos capazes de interagir com o mundo físico. Robôs móveis, manipuladores, veículos autônomos, drones e plataformas exploratórias passaram a ser empregados em atividades que exigem percepção do ambiente, tomada de decisão e execução de movimentos seguros. Com isso, a robótica se consolidou como uma área interdisciplinar, envolvendo conceitos de computação, controle, inteligência artificial, mecânica, eletrônica e sistemas embarcados. %Tem mais algum de robo humanoide com isso que representa usa rrt e como posso colocar aqui?

Sob esse ponto, a navegação autônoma assume papel fundamental porque para que um robô possa cumprir uma tarefa, não basta apenas possuir sensores e atuadores, também é necessário que ele seja capaz de determinar uma sequência de movimentos que o conduza de uma configuração inicial até uma configuração objetivo, evitando colisões e respeitando as restrições impostas pelo ambiente e pelo próprio sistema robótico. Esse problema é conhecido como planejamento de caminho, sendo uma das questões clássicas da robótica e da computação aplicada \citep{lavalle2006planning}.

Embora possa ser descrito de maneira simples, o planejamento de caminho apresenta dificuldades relevantes. O espaço no qual o robô se move pode conter obstáculos, regiões estreitas, áreas desconectadas ou caminhos de difícil acesso. Além disso, quando o robô possui muitos graus de liberdade, como ocorre em manipuladores, o espaço de busca cresce rapidamente, tornando inviável a análise exaustiva de todas as possibilidades. Em robôs móveis, mesmo em ambientes bidimensionais, a presença de obstáculos, restrições de movimento e necessidade de resposta em tempo reduzido já tornam o problema desafiador.

Tradicionalmente, muitos problemas de navegação podem ser tratados por métodos baseados em grades, contudo essas abordagens são eficientes apenas em cenários nos quais o espaço de busca pode ser discretizado de maneira adequada. Em ambientes contínuos, de alta dimensionalidade ou com grande quantidade de restrições geométricas, a discretização completa do espaço pode se tornar computacionalmente custosa. Além disso, a escolha da resolução da grade influencia diretamente a qualidade da solução: resoluções muito baixas podem ocultar passagens estreitas, enquanto resoluções muito altas aumentam significativamente o custo computacional da busca.


%OMPL is the default planning library in MoveIt and has been used for many robots. Here, OMPL is used to plan footsteps for NASA's Robonaut2 aboard the International Space Station

Para contornar essa limitação, surgiram os algoritmos de planejamento baseados em amostragem, que constroem uma representação parcial do espaço a partir de amostras aleatórias, conectando regiões livres na forma de grafos ou árvores. Essa abordagem lida naturalmente com espaços contínuos e de alta dimensionalidade \cite{lavalle1998rapidly}, e está consolidada em bibliotecas amplamente adotadas, como a Open Motion Planning Library (OMPL), que disponibiliza algoritmos como PRM e RRT, além de integração com o MoveIt/ROS e o simulador CoppeliaSim \cite{ompl,moveit_ompl}. A própria documentação da OMPL menciona seu uso no planejamento de movimentos do Robonaut2 da NASA a bordo da Estação Espacial Internacional, ilustrando a relevância prática desses métodos \cite{ompl}.

Diante desse cenário, este trabalho tem como foco o estudo e a análise de algoritmos de planejamento de movimento baseados em amostragem, com ênfase nos algoritmos RRT e EST. Busca-se compreender como essas abordagens exploram o espaço de busca, como se comportam diante de obstáculos e passagens restritas, e quais características influenciam sua capacidade de encontrar caminhos viáveis.

%\lipsum[1-4]

\section{Objetivos}

Objetivo geral

\subsection{Objetivos Específicos}
\begin{itemize}
\item Objetivo Específico 1

\item Objetivo Específico 2

\item Objetivo Específico 3

\end{itemize}

\section{Organização do trabalho}

\lipsum[1]